A permitted seating is a permutation of the eight chairs in which each nonfixed person moves one step around the circle. Any cycle of length two swaps adjacent people. A cycle of length at least three must continue in the same direction: changing direction would immediately revisit the previous chair. It therefore uses the entire circle. There are two such full rotations.

All other permutations are disjoint swaps along edges of the cycle graph. The number of sets of $k$ nonadjacent edges in a cycle of length $n$ is $\dfrac{n}{n-k}\binom{n-k}{k}$: count line matchings according to whether the wraparound edge is selected, or simplify $\binom{n-k}{k}+\binom{n-k-1}{k-1}$. For $n=8$ and $k=0,1,2,3,4$, the counts are $1,8,20,16,2$. Thus
\[
\boxed{1+8+20+16+2+2=49}.
\]
\begin{asy}
texpreamble("\usepackage{amsmath,amssymb}");
size(240);
defaultpen(fontsize(9pt));
for(int i=0;i<8;++i){pair p=2dir(90-45i);draw(p--2dir(90-45(i+1)),gray);dot(p);label(string(i+1),2.3dir(90-45i));}draw(2dir(90)--2dir(45),blue+1.5bp);draw(2dir(-45)--2dir(-90),blue+1.5bp);label("Two disjoint swaps",(0,-3));
\end{asy}
